\section{Labels, Outcomes and Policy Context}\label{sec:target-audit}
The extraction distinguishes the label being predicted from the action recorded, the original task's outcome and its policy context. The forty-four rows in \path{field-map/target-evidence.csv} retain the source version, setting, passage locator and each field's status. \emph{Reported} means explicit evidence in the inspected passage; \emph{partial} records a supplied component and a specific gap; \emph{not located} applies to the inspected scope; \emph{not applicable} identifies a field outside the endpoint. These statuses describe evidence availability, not study quality.

Three contrasts illustrate the rules. WRA's prompt target is an assigned research query, whereas ANT-stage uses log-derived primitive labels; recognizing the first does not establish the second. Trace excludes forty of 1,050 planned clean runs, while AgentTrace uses parser outcomes and ordered repairs as probe signals. These are session eligibility and within-session behavior, respectively, not a paired estimate of the effect of failure filtering. ANT-risk labels benchmark misuse; AgentTell evaluates action choice under a disclosure rule, and CHeaT reports a defined CTF-prevention outcome. None supplies a general per-user authorization label for arbitrary executions.

Each contrast remains tied to its sources. The source-withdrawal analysis in Section~\ref{supp:conclusion-withdrawal} states which required premises survive each removal. The complete extraction is provided as CSV in the review artifact, alongside the full earlier typeset table in \path{historical-studies/r98-supplement-archive.pdf}. Labels, outcomes and policy fields have not been inferred from classifier accuracy or completed from missing source information.
